% Part: normal-modal-logic
% Chapter: completeness
% Section: lindenbaums-lemma

\documentclass[../../../include/open-logic-section]{subfiles}

\begin{document}

\olfileid[zh]{nml}{com}{lin}

\olsection{Lindenbaum 引理}

Lindenbaum 引理断言，每一组 $\Sigma$-一致的!!{formula}集都包含在至少一个 \emph{完全的}
$\Sigma$-一致的集合中。我们对典范模型的构造将表明，对每个完全的 $\Sigma$-一致的集合~$\Delta$，
在典范模型中都存在一个世界，使得所有且仅有~$\Delta$ 中的!!{formula}在该世界处为真。
因此，Lindenbaum 引理保证了每个 $\Sigma$-一致的集合在典范模型的某个世界处为真。

\begin{thm}[Lindenbaum 引理]\ollabel{thm:lindenbaum}
  如果 $\Gamma$ 是 $\Sigma$-一致的，则存在一个完全的
  $\Sigma$-一致的集合 $\Delta$ 扩张~$\Gamma$。
\end{thm}

\begin{proof}
设 $!A_0$、$!A_1$、\dots{} 是该语言所有公式的一个穷尽性列出（允许重复）。例如，
先列出 $\Obj p_0$，并在每个阶段 $n \ge 1$ 列出仅使用取自 $\Obj p_0$、\dots,~$\Obj p_n$
之中的变元的长度为~$n$ 的有穷多个公式。我们通过关于~$n$ 的归纳来定义!!{formula}的集合
$\Delta_n$，然后置 $\Delta = \bigcup_n \Delta_n$。我们先令 $\Delta_0 = \Gamma$。
假设 $\Delta_n$ 已被定义，我们如下定义 $\Delta_{n+1}$：
\[
\Delta_{n+1} =
\begin{cases}
  \Delta_n \cup \{!A_n\}, & \text{如果 $\Delta_n \cup \{ !A_n\}$
    $\Sigma$-一致；}\\
  \Delta_n \cup \{ \lnot !A_n\}, & \text{否则。}
\end{cases}
\]
现在设 $\Delta = \bigcup_{n=0}^\infty \Delta_n$。

我们必须表明，这个定义实际上给出一个具有所需性质的集合 $\Delta$，即 $\Gamma \subseteq \Delta$，
且 $\Delta$ 是完全 $\Sigma$-一致的。

显然 $\Gamma \subseteq \Delta$，因为由构造 $\Delta_0 \subseteq \Delta$，且 $\Delta_0 = \Gamma$。
事实上，对所有~$n$，$\Delta_n \subseteq \Delta$，因为 $\Delta$ 是所有~$\Delta_n$ 的并。
（由于在构造的每一步中，我们都往已经构造的集合中加入!!a{formula}，故 $\Delta_n \subseteq
\Delta_{n+1}$，从而由于 $\subseteq$ 是传递的，只要 $n \le m$，就有 $\Delta_n \subseteq
\Delta_{m}$。）在构造的每一阶段，我们要么加入 $!A_n$ 要么加入 $\lnot !A_n$，而每个!!{formula}都（至少一次）出现在所有~$!A_n$ 的列表中。因此，对每个 $!A$，要么 $!A \in \Delta$，
要么 $\lnot !A \in \Delta$，所以按定义 $\Delta$ 是完全的。

最后，我们必须表明 $\Delta$ 是 $\Sigma$-一致的。为此，我们证明（a）如果 $\Delta$ 是
$\Sigma$-不一致的，则某个 $\Delta_n$ 会是 $\Sigma$-不一致的，以及（b）所有 $\Delta_n$
都是 $\Sigma$-一致的。

于是假设 $\Delta$ 是 $\Sigma$-不一致的。则 $\Delta \Proves[\Sigma] \lfalse$，即存在
$!A_1$、\dots,~$!A_k \in \Delta$ 使得 $\Sigma \Proves !A_1 \lif (!A_2 \lif \cdots (!A_k
\lif \lfalse)\dots)$。由于 $\Delta = \bigcup_{n=0}^\infty \Delta_n$，每个 $!A_i \in \Delta_{n_i}$
对某个~$n_i$ 成立。设 $n$ 为这些数中最大的。由于 $n_i \le n$，$\Delta_{n_i} \subseteq \Delta_n$。
所以，所有 $!A_i$ 都在某个~$\Delta_n$ 中。这意味着 $\Delta_n \Proves[\Sigma] \lfalse$，
即 $\Delta_n$ 是 $\Sigma$-不一致的。

为了证明每个 $\Delta_n$ 都是 $\Sigma$-一致的，我们使用关于~$n$ 的简单归纳。
$\Delta_0 = \Gamma$，而我们假设 $\Gamma$ 是 $\Sigma$-一致的。所以该断言在 $n = 0$ 时成立。
现在假设它对 $n$ 成立，即 $\Delta_n$ 是 $\Sigma$-一致的。若 $\Delta_n \cup \{!A_n\}$ 是
$\Sigma$-一致的，则 $\Delta_{n+1}$ 要么是 $\Delta_n \cup \{!A_n\}$，否则是
$\Delta_n \cup \{\lnot!A_n\}$。在第一种情形，$\Delta_{n+1}$ 显然是 $\Sigma$-一致的。
然而，由
\olref[prf][con]{prop:consistencyfacts}\olref[prf][con]{prop:consistencyfacts-c}，
$\Delta_n \cup \{!A_n\}$ 或 $\Delta_n \cup \{\lnot!A_n\}$ 之一是
一致的，所以在另一种情形下 $\Delta_{n+1}$ 也是一致的。
\end{proof}

\begin{cor}\ollabel{cor:provability-characterization}
  $\Gamma \Proves[\Sigma] !A$ 当且仅当对每个 扩张 $\Gamma$ 的完全 $\Sigma$-一致的集合
  $\Delta$，有 $!A \in \Delta$
  （包括 $\Gamma = \emptyset$ 的情形，此时我们得到模态系统 $\Sigma$ 的另一种刻画。）
\end{cor}

\begin{proof}
  设 $\Gamma \Proves[\Sigma] !A$，并设 $\Delta$ 是任意扩张 $\Gamma$ 的完全 $\Sigma$-一致的集合。
  如果 $!A \notin \Delta$，则由极大性 $\lnot!A \in \Delta$，因此
  $\Delta \Proves[\Sigma] !A$（由单调性）且 $\Delta \Proves[\Sigma] \lnot!A$
  （由自反性），从而 $\Delta$ 是不一致的。反之，如果 $\Gamma \Proves/[\Sigma] !A$，则
  $\Gamma \cup \{ \lnot!A\}$ 是 $\Sigma$-一致的，且由 Lindenbaum 引理存在一个扩张
  $\Gamma \cup \{ \lnot!A \}$ 的完全一致集合 $\Delta$。由一致性，$!A \notin \Delta$。
\end{proof}

\end{document}
